\section{Limitations \& Future Prospects}
\label{sec:future-work}

Although our results provide initial evidence for the feasibility, competitiveness, and scaling potential of \method, we view the current system as a starting point rather than a finished endpoint. First, the model remains far from convergence, especially the latent DiT, while the present training scale is limited by available resources. Larger models, longer training, and broader data and compute budgets are therefore needed to establish its performance ceiling. Second, compression reduces the dominant linear and attention FLOPs of the latent prior in theory, but end-to-end generation speed also depends on the decoder, CFG branches, latent feedback, and hardware utilization. Infrastructure co-design---including specialized kernels, cache management, and parallel scheduling---is needed to translate theoretical savings into practical speedups. Third, \method already shows compatibility with autoregressive and masked-diffusion decoders, but their co-training objectives, inference trade-offs, and optimal division of labor remain underexplored. Finally, the continuous latent space may provide a shared interface between text and other continuous modalities such as vision, opening a path toward unified latent representations, cross-modal priors, and unified understanding and generation.
